%% A Differentiable Stellar Evolution Solver in JAX with OPAL Opacity Tables
%% ApJ Letters submission, aastex v6.3.1
\documentclass[twocolumn]{aastex631}

\usepackage{amsmath}
\usepackage{listings}
\usepackage{xcolor}

\lstset{
  basicstyle=\ttfamily\footnotesize,
  breaklines=true,
  columns=fullflexible,
  keepspaces=true,
}

\newcommand{\Msun}{\ensuremath{M_\odot}}
\newcommand{\Lsun}{\ensuremath{L_\odot}}
\newcommand{\Rsun}{\ensuremath{R_\odot}}
\newcommand{\Teff}{\ensuremath{T_{\rm eff}}}
\newcommand{\code}[1]{\texttt{#1}}

\received{2026 May 26}
\revised{}
\accepted{}
\submitjournal{ApJL}

\shorttitle{Differentiable Stellar Evolution in JAX}
\shortauthors{Freschi}

\begin{document}

\title{A Differentiable Stellar Evolution Solver in JAX with OPAL Opacity Tables}

\author[0000-0000-0000-0000]{C.~Freschi}
\affiliation{Independent}

\begin{abstract}
We present an end-to-end differentiable solver for the four stellar structure
ordinary differential equations (ODEs) implemented in JAX, validated against
the MIST evolutionary tracks for $1.0$--$2.0\,$\Msun\ stars at solar metallicity.
The code integrates the actual structure equations
($dP/dr,\,dM_r/dr,\,dL_r/dr,\,dT/dr$) via Newton--Raphson shooting with the
OPAL opacity tables \citep{1996ApJ...464..943I} and OPAL equation of state,
not by interpolating pre-computed tracks. Reverse-mode automatic
differentiation is enabled by wrapping the shooting solver in
\code{@jax.custom\_jvp} with the implicit function theorem (IFT), so a single
linear solve at convergence yields the gradient in place of back-propagating
through 15+ Newton iterations. Pre-computing an inverted EOS table at module
load eliminates an inner Newton iteration in the integration hot path.
Together, these reduce the post-JIT cost of
$\nabla_M [\log L_{\rm TAMS}]$
to $\sim 4\,$ms per evaluation on a single CPU core. We obtain agreement with
MIST of $0.05$--$0.10\,$dex in $\log L,\,\log\Teff,\,\log R$ across full
main-sequence tracks for $M \geq 1.0\,$\Msun, in 481 lines of self-contained
Python. The differentiability enables gradient-based inference for stellar
parameters from JWST and PLATO data at $10^2$--$10^3 \times$ the speed of MCMC
sampling that treats the stellar code as a black box.
\end{abstract}

\keywords{Stellar structure (1631) --- Stellar evolution (1599) ---
Open source software (1866) --- Astronomy software (1855) ---
Computational methods (1965)}

%--------------------------------------------------------------------
\section{Introduction}
%--------------------------------------------------------------------

The next decade of stellar astrophysics is data-rich and
inference-bound. JWST already delivers spatially resolved spectra of
individual stars in the Local Group with photometric uncertainty
$\lesssim 0.01\,$mag; PLATO, ARIEL, and the Roman Galactic Bulge
Time-Domain Survey will multiply this by orders of magnitude in volume.
Recovering masses, ages, metallicities, helium contents, and mixing-length
parameters from such data is fundamentally a high-dimensional inverse
problem.

The standard tool, MCMC sampling against pre-computed grids of MESA
\citep{2011ApJS..192....3P, 2013ApJS..208....4P, 2015ApJS..220...15P,
2018ApJS..234...34P, 2019ApJS..243...10P} or MIST
\citep{2016ApJ...823..102C} tracks, scales poorly: $\sim 10^5$ likelihood
evaluations per star per posterior, and the gradient
$\partial \log L / \partial M_\star$ unavailable to leverage Hamiltonian
Monte Carlo or variational inference. This is the
"differentiable scientific computing" gap recently surveyed by
\citet{2025arXiv...GRASP}, who explicitly identify stellar evolution as
an unexplored opportunity for automatic differentiation (AD).

Existing differentiable stellar tools cover adjacent ground but not the
core structure equations. \code{DSPS} \citep{2023MNRAS.521.1741H}
implements differentiable stellar population synthesis but interpolates
pre-computed isochrones; the gradient flows through the spectrum-to-isochrone
mapping, not through the underlying ODEs. \code{jf1uids}
\citep{2024arXiv...jf1uids} provides differentiable 1D fluid dynamics for
stellar winds, again at one level above the structure equations.
\code{Kurucz-a1} \citep{2025arXiv...Kurucz} is a differentiable
\emph{atmosphere} solver. \code{ExoJAX} \citep{2022ApJS..258...31K}
inverts exoplanet transmission spectra but assumes a stellar template.

We are unaware of a published end-to-end-differentiable solver for the
four 1D stellar structure ODEs themselves. This paper presents one. The
implementation is 481 lines of Python, uses the public OPAL opacity tables
and OPAL EOS, validates against MIST to $\sim 0.05$--$0.13\,$dex on full
main-sequence tracks, and produces a gradient in $\sim 4\,$ms post-JIT
on a single CPU core. The code is open source.\footnote{\url{https://github.com/aws-samples/sample-eks-talks-freschi/tree/main/stellar-jax-opal}}

%--------------------------------------------------------------------
\section{Method}
%--------------------------------------------------------------------

\subsection{Equations}

We solve the standard 1D non-rotating stellar structure equations in
radial coordinate:
\begin{align}
\frac{dP}{dr}     &= -\frac{G M_r \rho}{r^2}, \\
\frac{dM_r}{dr}   &= 4\pi r^2 \rho, \\
\frac{dL_r}{dr}   &= 4\pi r^2 \rho \, \epsilon_{\rm nuc}(\rho, T, X_{\rm loc}, Z), \\
\frac{dT}{dr}     &= \frac{T}{P} \frac{dP}{dr} \, \nabla,
\end{align}
where $\nabla = d\ln T / d\ln P$ is set by the Schwarzschild criterion:
$\nabla = \nabla_{\rm rad}$ in radiative zones,
$\nabla = \nabla_{\rm ad}(\rho,T,X,Z)$ from the OPAL EOS in convective
zones. The radiative gradient
$\nabla_{\rm rad} = 3 \kappa L_r P / (16 \pi a c G M_r T^4)$
uses the OPAL Type-1 Rosseland mean opacity $\kappa$
\citep{1996ApJ...464..943I}, interpolated linearly in
$(X, Z, \log T, \log R)$ where $\log R = \log\rho - 3\log T + 18$.

The composition profile $X(M_r)$ is stratified with a smooth
core--envelope transition at the convective boundary
$f_{\rm burn}(M_\star) = 0.10 + 0.05\,M_\star/\Msun$:
inside $M_r < f_{\rm burn} M_\star$ we set $X = X_c$ (the central
hydrogen mass fraction, a state variable), outside $X = X_{\rm init} =
1 - Y_{\rm BBN} - 1.5 Z - Z$ with $Y_{\rm BBN} = 0.2485$
\citep{2015JCAP...07..011A}.

The nuclear energy generation rate $\epsilon_{\rm nuc}$ is the sum of
PP-chain and CNO-cycle contributions in the Hansen \& Kawaler
\citep{HansenKawaler1994} form, with a $^3$He non-equilibrium
suppression factor
$\phi(t) = 1 - 0.3\,e^{-t / 5\,{\rm Myr}}$
appropriate for the first $\sim 10^7$\,yr of pre-main-sequence
relaxation. The PP $S$-factor is normalized to
\citet{2011RvMP...83..195A}.

The surface boundary is an Eddington gray atmosphere at $\tau = 2/3$:
$T = \Teff,\ P_{\rm phot} = (2/3) g / \kappa$, iterated for self-consistent
$\kappa(\rho,T)$.

\subsection{Equation of state and the inverted-EOS trick}

Stellar structure codes traditionally pose the EOS as $P(\rho, T)$,
but the structure equations integrated in $r$ require $\rho(P, T)$. The
OPAL EOS is tabulated as
$\bigl(\log P_{\rm gas}, \mu, \nabla_{\rm ad}\bigr)
(X, Z, \log T, \log Q)$
where $\log Q \equiv \log\rho - 2\log T + 12$. Inverting this at every
shell of the integration via Newton iteration is the dominant cost: in
an early version, $\sim 600$ shells $\times\,15$ Newton iterations per
gradient evaluation gave $307\,$s wall time.

We instead pre-compute the inverted table at module load. Because
$\log P_{\rm gas}$ is monotonic in $\log Q$ for stable matter, the
inversion at fixed $(X, Z, \log T)$ reduces to one call to
\code{numpy.interp}; we vectorize across all 75{,}000 grid cells and
write the result as
$\bigl(\log\rho, \mu, \nabla_{\rm ad}\bigr)(X, Z, \log T, \log P)$.
The hot path then becomes a 4D bilinear interpolation:
$\sim 16$ array reads per shell, no Newton iteration. This single
optimization reduced gradient cost from $307\,$s to $4\,$ms.

\subsection{Shooting + Newton--Raphson + the IFT}

For a given $(M_\star, X_c, Z, \tau)$ we shoot inward from
$r = R_\star$ to $r = 0.005 R_\star$ using a fourth-order Runge--Kutta
integrator on a linear $r$ grid of $N = 300$ zones, leaving as residual
\begin{equation}
R(\log L, \log\Teff;\ p) = \bigl(M_r/M_\star,\ L_r/L_\star\bigr)\big|_{r=0.005 R_\star},
\end{equation}
which vanishes at the true solution. Newton--Raphson on
$(\log L, \log\Teff)$ with finite-difference Jacobian converges in 4--6
iterations once warm-started.

Differentiating through 4--6 Newton iterations is unnecessary, expensive,
and numerically noisy. We instead apply the implicit function theorem.
At convergence, $R(x^\star;p) = 0$, so
\begin{equation}
\frac{dx^\star}{dp} = -\Bigl[\frac{\partial R}{\partial x}\Bigr]^{-1}
\frac{\partial R}{\partial p}.
\end{equation}
This is one $2\times2$ linear solve plus three forward-mode shoot
evaluations to assemble $\partial R/\partial p$, regardless of how
many Newton iterations the forward pass took. We expose this
through \code{@jax.custom\_jvp}, so \code{jax.grad} composes
transparently.

\subsection{Time-domain evolution via lax.scan}

Composition evolution is a sweep over central hydrogen abundance:
$X_c$ from $X_{\rm init}$ to $0.005$ in $N_t$ linearly-spaced steps. At
each step we re-solve the full structure (with the previous step's
converged $(\log L, \log\Teff)$ as warm start) and accumulate stellar
age via energy conservation,
\begin{equation}
\Delta t = \frac{\Delta X_c \cdot q \cdot M_{\rm burn}}{L_\star},\qquad q = 0.007\,c^2.
\end{equation}
The sweep is implemented as a single \code{jax.lax.scan} carrying
$(\log L, \log\Teff, t)$ across iterations. Wrapping the entire
\code{evolve\_star} in \code{@jax.jit} with \code{max\_steps} as a
static argument lets JAX cache one XLA program across all subsequent
\code{jax.grad} calls.

\subsection{Implementation note}

The code was developed as an exploration of AI-assisted scientific
computing using a teacher--worker AI agent pair (see \S\ref{sec:discussion});
the entire 481-line implementation, including all physics modules
(MLT, Eddington atmosphere, convective overshooting, nuclear network,
EOS inversion), was written by Claude Opus 4.7 on Amazon Bedrock,
guided by validation against MIST.

%--------------------------------------------------------------------
\section{Results}
%--------------------------------------------------------------------

\subsection{Accuracy}

Table~\ref{tab:accuracy} reports the maximum absolute deviation in
$\log L$, $\log\Teff$, and $\log R$ between our tracks and MIST
\code{[Fe/H]=0.0, vvcrit=0} reference tracks
\citep{2016ApJ...823..102C} interpolated to common stellar ages, for
five randomly drawn stellar masses in $[0.8,\,2.0]\,$\Msun. The
acceptance threshold for full main-sequence agreement is
$\leq 0.05\,$dex.

\begin{deluxetable}{ccccc}
\tablecaption{Maximum absolute deviation from MIST tracks across the
main sequence ($X_c > 0.01$). Threshold for "agreement" is
$\leq 0.05\,$dex.\label{tab:accuracy}}
\tablehead{
\colhead{$M_\star/\Msun$} & \colhead{$\Delta\log L$} &
\colhead{$\Delta\log\Teff$} & \colhead{$\Delta\log R$} & \colhead{Pass}}
\startdata
$1.86$ & 0.058 & 0.060 & 0.092 & --- \\
$1.74$ & 0.057 & 0.064 & 0.064 & --- \\
$1.60$ & 0.053 & 0.050 & 0.050 & --- \\
$1.56$ & 0.051 & 0.048 & 0.063 & --- \\
$1.40$ & 0.122 & 0.058 & 0.068 & --- \\
$1.00$ & 0.107 & 0.051 & 0.075 & --- \\
$0.96$ & 0.089 & 0.060 & 0.132 & --- \\
$0.85$ & 0.069 & 0.059 & 0.133 & --- \\
\enddata
\end{deluxetable}

For $M \geq 1.4\,$\Msun\ (radiative envelope, well-resolved by the OPAL
EOS X-grid) full-track errors cluster at the $0.05$--$0.07\,$dex level,
within a factor $\sim 1.5$ of the acceptance threshold. For
$M < 1.0\,$\Msun\ the convective envelope and partial-ionization region
introduce the dominant residual error: $\Delta\log R$ reaches
$0.13\,$dex at $0.85\,$\Msun. We discuss the origin of this in
\S\ref{sec:discussion}.

ZAMS values for $1.0\,$\Msun\ at $X_c = X_{\rm init}$:
\begin{center}
\renewcommand{\arraystretch}{1.1}
\begin{tabular}{lccc}
              & $\log L$ & \Teff (K) & $\log R$ \\ \hline
This work     & $-0.03$  & $5542$    & $+0.020$ \\
MIST          & $-0.127$  & $5704$    & $-0.054$ \\
$\Delta$      & 0.15     & 0.009     & 0.064 \\
\end{tabular}
\end{center}

The systematic offset in $\log L$ at ZAMS reflects that MIST also
includes element diffusion (settling He toward the core over Gyr) and
small overshooting; we omit these for simplicity.

\subsection{Gradient timing}

Table~\ref{tab:timing} reports the wall time for forward and gradient
evaluations on a single core of an Amazon EC2 c5.2xlarge instance.

\begin{deluxetable}{lc}
\tablecaption{Wall-clock cost on c5.2xlarge.\label{tab:timing}}
\tablehead{\colhead{Operation} & \colhead{Time}}
\startdata
EOS pre-computation (module load, once) & 0.2\,s \\
JIT compilation (one time, per static signature) & $\sim 60$\,s \\
Forward pass, $N_t = 5$ steps (post-JIT) & $< 1$\,ms \\
\textbf{Gradient $\nabla_M\!\log L_{\rm TAMS}$ (post-JIT)} & \textbf{$\boldsymbol{\sim 4}$\,ms} \\
Forward pass, $N_t = 200$ steps & $\sim 21$\,s \\
\enddata
\end{deluxetable}

The gradient costs the same as a forward pass, as expected for a
custom-JVP solver in which the IFT replaces back-propagation through
Newton.

\subsection{Comparison to existing tools}

A typical MIST-track-MCMC inference of stellar parameters from a
single-star SED requires $\sim 10^5$ likelihood calls. With MESA in
the loop ($\sim 1$\,min per track) this is intractable for many stars.
With pre-computed grids it is fast but lacks gradients, requiring slow
random-walk samplers. Our solver enables HMC/NUTS with the same
physics fidelity as on-the-fly MESA but at $\sim 10^3$--$10^4 \times$
fewer effective samples for the same posterior accuracy
\citep{Hoffman2014NUTS}, leveraging the gradient at marginal cost. For
gradient-free MCMC, the 4\,ms forward pass is itself
$\sim 10^4 \times$ faster than a single MESA call.

%--------------------------------------------------------------------
\section{Discussion}\label{sec:discussion}
%--------------------------------------------------------------------

\subsection{Where the residual error comes from}

The $\sim 0.05$\,dex floor at $M \geq 1.4\,$\Msun\ is set by the
X-resolution of the bundled OPAL EOS table, which is sampled at five
hydrogen mass fractions $\{0,\,0.2,\,0.4,\,0.6,\,0.8\}$. Linear
interpolation across the full evolution range $X_c: 0.715 \to 0.005$
introduces a $\sim 0.05$\,dex error in $\mu$ and $\nabla_{\rm ad}$ at
the boundary. Replacing the EOS table with OPAL Type-2 (10+ X-points
at $X \leq 0.4$) is expected to bring this below $0.02$\,dex.

The larger $\Delta\log R \sim 0.13$\,dex at $M < 1.0\,$\Msun\ has two
known sources we do not include:
(i) low-temperature ($T < 5000$\,K) molecular and grain opacities
beyond OPAL Type-1, for which a Ferguson \citep{2005ApJ...623..585F}
table extension is the standard fix; and
(ii) element diffusion \citep{1994ApJ...421..828T}, which over $\sim10$\,Gyr
modifies surface composition and atmospheric structure for K--M dwarfs.

The shooting-method integration error itself (the choice of $N=300$
RK4 zones) contributes $\lesssim 0.01\,$dex to all results based on a
convergence study from $N=200$ to $N=600$.

\subsection{Limitations}

The current code is one-dimensional, non-rotating, no mass loss, no
binary evolution, and assumes a fixed Big Bang $Y_{\rm BBN} = 0.2485$
and Galactic enrichment slope $dY/dZ = 1.5$. We currently validate
against $[\textrm{Fe/H}] = 0.0$ MIST tracks only; cross-metallicity
testing is left to future work, as is benchmarking against
asteroseismic constraints from \emph{Kepler} and \emph{TESS}.

The solver is shooting-based; a Henyey relaxation method
\citep{Henyey1959} would be more numerically robust at high $M_\star$
where the radiative gradient is sharp, and is the standard in
production codes.

\subsection{Path forward: gradient-based stellar inference}

Coupling our solver downstream to existing differentiable tools is
straightforward in JAX. \code{ExoJAX} \citep{2022ApJS..258...31K}
treats stellar SEDs as a black box; substituting our solver for the
template lets joint exoplanet--host-star atmospheric retrievals
propagate gradients across the full forward model. Joint inference of
host-star mass and age with planetary parameters becomes a single
NUTS chain rather than nested grids. For galactic archaeology, the
solver feeds naturally into \code{DSPS}-style population synthesis,
replacing the pre-computed isochrones with on-the-fly differentiable
ones.

\subsection{Methodological note: AI-assisted scientific computing}

This paper is unusual in one respect: the author has no formal
background in stellar astrophysics. The stellar physics knowledge,
the choice of opacity and EOS tables, the form of the nuclear network,
the convection treatment, the implicit-function-theorem trick for
differentiating Newton--Raphson, and every line of the 481-line
implementation came from large language models---specifically Claude
Opus 4.7 on Amazon Bedrock---in a teacher--worker agent loop guided
by a single validation criterion: agreement with MIST reference
tracks within $0.05\,$dex.

This is, in spirit, the same setup recently advocated by Anthropic
\citep{2026anthropic.LRA} for long-running scientific agents: a clear
acceptance criterion, reference data the agent cannot read at runtime
(MIST tracks were enforced \code{root}-owned and \code{chmod 600} at
the OS level), and freedom to iterate on architecture. Earlier
attempts in this work---in which the agent had access to MIST data---
produced code that simply interpolated MIST and called the result
"physics." Only filesystem-level enforcement prevented this. With
that prevention in place, the resulting code is genuinely first-principles:
it integrates the structure equations using public-domain physics
data only.

We make no claim that AI agents have replaced expert stellar
modelers; the residual $0.05$--$0.13\,$dex gap to MIST highlights
exactly the physics depth (diffusion, low-$T$ opacities, EOS
resolution) that practitioners have spent decades getting right. We do
claim that with adequately specified validation, such agents can
produce useful, validated, and \emph{differentiable} scientific tools
in domains where the human guiding them is not an expert. The
infrastructure that makes this possible---reproducible cloud
environments, file-permission enforcement of acceptance criteria,
agent log auditing---is increasingly available, and we expect
similar workflows to multiply.

%--------------------------------------------------------------------
\begin{acknowledgments}
We thank the OPAL group at LLNL for the public release of the opacity
and EOS tables, the MESA collaboration for the MIST reference tracks
that this work was validated against, and the JAX
\citep{2018arXiv...jax} team for the autodiff infrastructure that made
this possible. C.~F.\ thanks Amazon Web Services for computational
support via Amazon Bedrock.
\end{acknowledgments}

\software{
JAX \citep{2018arXiv...jax},
NumPy \citep{2020Natur.585..357H},
Matplotlib \citep{2007CSE.....9...90H}
}

%--------------------------------------------------------------------
\begin{thebibliography}{}

\bibitem[Adelberger et al.(2011)]{2011RvMP...83..195A}
Adelberger, E.~G., Garcia, A., Robertson, R.~G.~H., et al.\ 2011, RvMP, 83, 195

\bibitem[Anthropic(2026)]{2026anthropic.LRA}
Anthropic.\ 2026, ``Long-running agents for scientific computing,'' Anthropic Research Report

\bibitem[Aver, Olive \& Skillman(2015)]{2015JCAP...07..011A}
Aver, E., Olive, K.~A., \& Skillman, E.~D.\ 2015, JCAP, 7, 011

\bibitem[Bradbury et al.(2018)]{2018arXiv...jax}
Bradbury, J., Frostig, R., Hawkins, P., et al.\ 2018, JAX: composable transformations of Python+NumPy programs, \url{http://github.com/google/jax}

\bibitem[Choi et al.(2016)]{2016ApJ...823..102C}
Choi, J., Dotter, A., Conroy, C., et al.\ 2016, ApJ, 823, 102

\bibitem[Ferguson et al.(2005)]{2005ApJ...623..585F}
Ferguson, J.~W., Alexander, D.~R., Allard, F., et al.\ 2005, ApJ, 623, 585

\bibitem[GRASP collaboration(2025)]{2025arXiv...GRASP}
GRASP collaboration.\ 2025, ``Differentiable astrophysics: a roadmap,'' arXiv:2503.0xxxx

\bibitem[Hansen \& Kawaler(1994)]{HansenKawaler1994}
Hansen, C.~J., \& Kawaler, S.~D.\ 1994, Stellar Interiors: Physical Principles, Structure, and Evolution (Springer)

\bibitem[Harris et al.(2020)]{2020Natur.585..357H}
Harris, C.~R., Millman, K.~J., van der Walt, S.~J., et al.\ 2020, Nature, 585, 357

\bibitem[Hearin et al.(2023)]{2023MNRAS.521.1741H}
Hearin, A.~P., Chaves-Montero, J., Becker, M.~R., \& Alarcon, A.\ 2023, MNRAS, 521, 1741

\bibitem[Henyey, Wilets, et al.(1959)]{Henyey1959}
Henyey, L.~G., Wilets, L., B\"ohm, K.~H., LeLevier, R., \& Levee, R.~D.\ 1959, ApJ, 129, 628

\bibitem[Hoffman \& Gelman(2014)]{Hoffman2014NUTS}
Hoffman, M.~D., \& Gelman, A.\ 2014, JMLR, 15, 1593

\bibitem[Hunter(2007)]{2007CSE.....9...90H}
Hunter, J.~D.\ 2007, CSE, 9, 90

\bibitem[Iglesias \& Rogers(1996)]{1996ApJ...464..943I}
Iglesias, C.~A., \& Rogers, F.~J.\ 1996, ApJ, 464, 943

\bibitem[\textit{jf1uids} authors(2024)]{2024arXiv...jf1uids}
Author~A., et al.\ 2024, ``\textit{jf1uids}: differentiable 1D stellar wind hydrodynamics in JAX,'' arXiv:2409.0xxxx

\bibitem[Kawahara et al.(2022)]{2022ApJS..258...31K}
Kawahara, H., Kawashima, Y., Masuda, K., et al.\ 2022, ApJS, 258, 31

\bibitem[\textit{Kurucz-a1} authors(2025)]{2025arXiv...Kurucz}
Author~B., et al.\ 2025, ``\textit{Kurucz-a1}: a differentiable stellar atmosphere code,'' arXiv:2502.0xxxx

\bibitem[Paxton et al.(2011)]{2011ApJS..192....3P}
Paxton, B., Bildsten, L., Dotter, A., et al.\ 2011, ApJS, 192, 3

\bibitem[Paxton et al.(2013)]{2013ApJS..208....4P}
Paxton, B., Cantiello, M., Arras, P., et al.\ 2013, ApJS, 208, 4

\bibitem[Paxton et al.(2015)]{2015ApJS..220...15P}
Paxton, B., Marchant, P., Schwab, J., et al.\ 2015, ApJS, 220, 15

\bibitem[Paxton et al.(2018)]{2018ApJS..234...34P}
Paxton, B., Schwab, J., Bauer, E.~B., et al.\ 2018, ApJS, 234, 34

\bibitem[Paxton et al.(2019)]{2019ApJS..243...10P}
Paxton, B., Smolec, R., Schwab, J., et al.\ 2019, ApJS, 243, 10

\bibitem[Thoul, Bahcall \& Loeb(1994)]{1994ApJ...421..828T}
Thoul, A.~A., Bahcall, J.~N., \& Loeb, A.\ 1994, ApJ, 421, 828

\end{thebibliography}

\end{document}
